\section{Introduction}
\label{sec:intro}

Suppose a data analyst filters a table of measurements with \texttt{WHERE -d > 'nan'::DOUBLE}. The database plans the
query, its optimizer rewrites the predicate into a cheaper form, and the analyst receives rows. Nothing fails and no
warning appears. If the rewrite was wrong, the analyst has no way to notice. Database testing tools such as
SQLancer~\cite{rigger2020tlp,rigger2020norec} and grammar-based fuzzers such as SQLsmith~\cite{sqlsmith} and
Squirrel~\cite{zhong2020squirrel} were built for this class of \emph{logic bug}. They generate many queries and compare
each answer with an answer that should be equal.

These tools work best on code that random queries reach often. An optimizer, however, is a long list of small rewrite
rules, and each rule fires only under narrow conditions: a constant of a specific type, a cast next to a comparison, a
join key of a list type. A rule that fires on one query in ten thousand is tested rarely, and its edge cases (type
limits, signed zero, NaN) more rarely still. A second problem is timing. Projects such as
DuckDB~\cite{raasveldt2019duckdb} are fuzzed continuously by their developers and by other researchers, so a bug found
in a release is often already fixed on the development branch. In our first month, every real bug our tools found in
DuckDB~1.5.6 had already been reported or fixed upstream (Section~\ref{sec:eval-methods}).

We focus on a third problem that starts where these tools stop: \emph{directing} the search at one optimizer rule at a
time, and proving that the generated tests actually exercise that rule. Large language models (LLMs) can read source
code and write programs, which makes directed generation possible~\cite{yang2024whitefox,xia2024fuzz4all}. But an LLM
also reports success it did not achieve. In an early run of our pipeline, 10 of 11 generators that an LLM agent wrote
and described as thorough produced exactly 8 queries per case and 5 to 12 distinct plans
(Section~\ref{sec:eval-judge}).

\paragraph{A real example.} Figure~\ref{fig:motivating} shows issue \#28052, which one of our LLM-written generators
found in DuckDB~2.0.0. DuckDB orders NaN above every other value, so \texttt{-d > NaN} is false for every row. The rule
\texttt{MoveUnaryMinusRule} rewrites \texttt{-x > c} into \texttt{x < -c}. Its guard handles a NaN \emph{column value}
but not a NaN \emph{constant}, so the query returns all 4 rows instead of none. With the expression rewriter switched
off, DuckDB returns the correct empty result, and DuckDB~1.5.6 was also correct. The maintainers merged a fix one day
after we filed (PR \#28071). This is not a single case: 8 of the 16 issues we filed are silent wrong results
(Table~\ref{tab:issues}).

\begin{figure}[t]
\begin{lstlisting}[basicstyle=\ttfamily\footnotesize,frame=single,columns=fullflexible]
CREATE TABLE f(d DOUBLE);
INSERT INTO f VALUES (1.5), (-2.5), ('nan'), ('inf');
SELECT d FROM f WHERE -d > 'nan'::DOUBLE ORDER BY d;
  -- DuckDB 2.0.0:  -2.5, 1.5, inf, nan     (4 rows, wrong)
  -- expected:      no rows                  (NaN is the largest value)
SET disabled_optimizers = 'expression_rewriter';
SELECT d FROM f WHERE -d > 'nan'::DOUBLE ORDER BY d;   -- no rows (correct)
\end{lstlisting}
\caption{Issue \#28052, found by an LLM-written generator. One rewrite rule turns an always-false predicate into an
always-true one. No error is raised; only a second run with the rule switched off shows the difference.}
\label{fig:motivating}
\end{figure}

Closest to our work, WhiteFox~\cite{yang2024whitefox} lets an LLM read compiler optimization code and write tests for
it, and directed greybox fuzzing~\cite{bohme2017aflgo} steers inputs toward chosen code locations. However, WhiteFox
targets compilers and trusts the LLM's description of what triggers an optimization, and directed greybox fuzzers need
distances over a control-flow graph, which is costly to compute for a large C++ database. Neither checks with the
database itself that a generated SQL test fired the intended rewrite.

\paragraph{Our approach.} We present \tool (Semantic LLM-Directed Fuzzing), a framework of four methods that run side
by side on the same infrastructure (Section~\ref{sec:approach}):
\begin{enumerate}
  \item \textbf{Method A, classic fuzzing}, as a baseline and safety net: SQLsmith on a sanitizer build, differential
        sweeps of DuckDB's own tests, and replay of past bug reports (Section~\ref{sec:method-a}).
  \item \textbf{Method B, feature-targeted fuzzers} for features that are new in DuckDB~2.0, such as triggers
        (Section~\ref{sec:method-b}).
  \item \textbf{Method C, new-code hints and LLM-written generators}: the LLM receives hints about code that is new or
        never executed, feature cards and fix dossiers, and writes a generator program (Section~\ref{sec:method-c}).
  \item \textbf{Method D, the rule-targeted loop}: a \emph{rule card} per optimizer rule, mined from its source and fix
        history, a hypothesis-first LLM generator, and an \emph{engine judge} that accepts the generator only if the rule
        fires on its queries (Section~\ref{sec:method-d}).
\end{enumerate}
All four methods feed one oracle layer and one triage pipeline. A person reads every surviving candidate before
anything is filed.

\paragraph{Results.} Over five weeks we filed 16 issues on DuckDB~2.0.0. Maintainers fixed 4 (\#27759, \#27850,
\#28052, \#28562), 9 more have open fix pull requests, and one of those (\#28547) is ours. Method D found 12 of the 16.
One unattended run of Method D spent 40 LLM calls on 13 targets, ran 32 generator rounds and turned 597 raw differences
into 9 triaged candidates and 6 issues (Table~\ref{tab:hunt1}). The main lesson is that the LLM is useful only when the
engine checks its work, and that the check must not be the only gate: four generators that found bugs failed our own
diversity floor, so we triage the hits of every generator but count only judged generators as test coverage.

\paragraph{Contributions.}
\begin{itemize}
  \item A framework of four complementary fuzzing methods for a database optimizer that share one judge, one oracle
        layer and one triage pipeline (Section~\ref{sec:approach}).
  \item Rule cards: per-rule targets mined from source code and fix history, including the early exits under which a
        rewrite fires and the places where it builds values (Section~\ref{sec:method-d}).
  \item An engine judge that measures rule-level reach and rejects generators that only claim it, checked by positive
        and negative controls (Sections~\ref{sec:judge} and~\ref{sec:eval-judge}).
  \item 16 upstream issues on DuckDB with minimal reproducers, 4 of them fixed, and one fix pull request
        (Section~\ref{sec:evaluation}). All code and data are public (Section~\ref{sec:data}).
\end{itemize}
